\title{FlashAttention-3: Fast and Accurate Attention with Asynchrony and Low-precision}

\begin{abstract}
Attention mechanisms, central to the Transformer architecture, constitute the principal computational constraint in both large language models and tasks requiring extended context handling. 
\textsc{FlashAttention} introduced a method to accelerate attention computations on GPUs by reducing unnecessary memory traffic. Nonetheless, it falls short of leveraging recent advances in hardware, with \textsc{FlashAttention-2} only achieving 35\% of the potential throughput on the H100 GPU. In this work, we propose three key optimizations to further accelerate attention on Hopper GPUs: utilizing asynchronous execution between Tensor Cores and TMA to (1) overlap computations and data transfers via warp specialization, (2) interleave block-level matrix multiplications with softmax steps, and (3) apply block-wise quantization and incoherent processing, capitalizing on enhanced hardware support for low-precision FP8 formats. Our approach, termed \textsc{FlashAttention-3}, demonstrates 1.5-2.0$\times$ speed improvements on H100 GPUs, achieving up to 740 TFLOPs/s in FP16 mode (corresponding to 75\% hardware utilization), and nearing 1.2 PFLOPs/s with FP8 precision. Empirically, we further establish that FP8 \textsc{FlashAttention-3} reduces numerical error by 2.6$\times$ when compared to a standard FP8 attention baseline.
\end{abstract}

\section{Introduction}
\label{sec:intro}

Within the Transformer framework~\citep{vaswani2017attention}, the attention component presents a central computational challenge due to its quadratic complexity with respect to sequence length in evaluating self-attention between queries and keys. Extending the reach of attention mechanisms to handle longer contexts promises to enable advancements such as understanding and reasoning across numerous lengthy documents~\citep{guo2021longt5,shaham2022scrolls,peng2023yarn} or files within vast code repositories~\citep{roziere2023code, li2023starcoder}. It also paves the way for processing novel data modalities, including high-resolution visual data~\citep{chen2022scaling}, audio~\citep{gulati2020conformer}, and video content~\citep{ho2022video}. Additionally, it supports emerging applications such as leveraging extensive user interaction histories~\citep{sun2019bert4rec} or managing agents’ operations over prolonged timelines~\citep{yao2022react}. Consequently, there is growing research aimed at accelerating attention computations for long contexts, which includes using approximate methods~\citep{katharopoulos2020transformers,choromanski2020rethinking,
tay2020efficient}, optimizing the underlying software (\citep{rabe2021self,dao2022flashattention,kwon2023efficient}), and even designing alternative model architectures~\citep{peng2023rwkv,sun2023retentive,gu2023mamba}.

This study extends the approach of \citet{dao2022flashattention}, who pioneered exact-attention algorithms that explicitly consider the GPU’s computational structure and hardware properties within their design. In their work, Dao et al. introduced \textsc{FlashAttention}, proposing an innovative tiling mechanism that allows attention computation to be parallelized while circumventing repeated access to slow global memory—this is achieved by consolidating all relevant operations within a single GPU kernel.

Building upon this, \citet{dao2023flashattention2} advanced the algorithm to create \textsc{FlashAttention-2}, which further distributes computation across the sequence length dimension and processes the forward pass’s inner loop on blocks of the key and value matrices. This redesign yields superior workload balancing and GPU occupancy. Nonetheless, our analysis reveals that \textsc{FlashAttention-2} still falls short in leveraging the full capabilities of state-of-the-art GPUs, with utilization rates around 35\%—substantially lower than the 80-90\% reached by highly tuned GEMM kernels on, for instance, the Hopper H100 GPU. One contributing factor is implementation nuances, such as the absence of Hopper-optimized instructions for Tensor Cores, with Ampere-era instructions used instead. Recent work, including ThunkerKitten~\citep{spector2024thunder} and cuDNN 9~\citep{cudnn9}, has demonstrated that the adoption of Hopper-specific instructions and tile-oriented programming models can deliver higher efficiency for attention mechanisms while streamlining their development.

At a deeper level, the algorithm behind \textsc{FlashAttention-2} is structured around a straightforward, synchronous paradigm and does not intentionally leverage either asynchrony or reduced-precision arithmetic within its framework. Asynchrony arises primarily due to the presence of specialized hardware components designed to optimize the most computationally demanding elements of machine learning tasks: for example, matrix multiplications executed on Tensor Cores or memory transfers carried out by Tensor Memory Accelerators (TMA), functioning independently from the main set of CUDA cores tasked with logical, integer, and floating-point calculations. The use of low-precision formats—such as FP8 introduced with Hopper and FP4 with Blackwell—continues the progression established by FP16 (from Pascal, 2017) and BF16 (from Ampere, 2020); this approach reliably delivers two to four times higher computational throughput given the same energy and silicon footprint. A detailed discussion of the Hopper hardware's support for these advancements can be found in \cref{subsec:hardware}. The core technical obstacle involves refactoring \textsc{FlashAttention-2} so it can exploit these modern hardware capabilities: utilizing asynchrony effectively means orchestrating simultaneous operations—such as matmul and softmax—even though these processes are inherently sequentially dependent; meanwhile, deploying low-precision techniques necessitates careful management of quantization artifacts, an issue made more critical by the occurrence of outlier features in large language models~\citep{dettmers2208llm, sun2024massive}.

To achieve this goal, we introduce \textsc{FlashAttention-3}, which integrates three novel techniques aimed at further enhancing efficiency on recent GPU architectures.\footnote{While our empirical findings are discussed using NVIDIA’s Hopper architecture as a reference, the proposed algorithm applies broadly to any GPU hardware that supports strong asynchronous execution and low-precision operations.}

\begin{enumerate}

\item \textbf{Producer-Consumer asynchrony:} We introduce a warp-specialized approach to software pipelining that leverages the asynchronous nature of data transfer and computation via Tensor Cores. This is accomplished by allocating data-producing and data-consuming tasks to distinct warps, which enhances the algorithm’s latency-masking capabilities for both memory accesses and instruction dispatch.

\item \textbf{Hiding softmax under asynchronous block-wise GEMMs:} We interleave the relatively less compute-intensive operations involved in softmax computation—such as floating point multiply-add and exponentiation—with the asynchronous execution of WGMMA GEMM instructions. This necessitates modifications to the \textsc{FlashAttention-2} procedure to alleviate certain sequential dependencies existing between the softmax and GEMM stages. Specifically, in the two-stage variant, softmax is carried out on a particular block of the score matrix concurrently as WGMMA, running asynchronously, computes subsequent blocks.

\item \textbf{Hardware-accelerated low-precision GEMM:} Our approach revises the forward pass algorithm to utilize FP8 Tensor Cores during GEMM operations, nearly achieving a twofold improvement in observed TFLOPs/s. This adaptation involves reconciling the divergent memory layouts required by WGMMA for blocks of FP32 accumulators and FP8 operand matrices. Techniques such as block quantization and incoherent processing are employed to counteract potential accuracy degradation stemming from the reduction to FP8 precision.
\end{enumerate}

In order to empirically assess the performance of our approach, we conduct benchmarks of \textsc{FlashAttention-3} on the H100 SXM5 GPU across a variety of settings. Our results indicate that (1) FP16 configuration achieves a 1.5-2.0$\times$ acceleration over \textsc{FlashAttention-2} in the forward computation (attaining up to 740 TFLOPs/s), and a 1.5-1.75$\times$ acceleration in the backward computation; (2) the FP8 format attains throughput near 1.2 PFLOPs/s; and (3) at high sequence lengths, FP16 demonstrates superior performance, with FP8 maintaining competitiveness\footnote{Specifically, for head dimension 64, the FP8 mode in \textsc{FlashAttention-3} has superior results, while for head dimensions 128 and 256, FP8 is comparable in non-causal masking contexts and lags when causal masking is present.} relative to the leading attention implementation available in NVIDIA’s cuDNN library. Additionally, we demonstrate that FP16 \textsc{FlashAttention-3} maintains the same level of numerical error as \textsc{FlashAttention-2} and provides improved accuracy compared to the conventional attention baseline, since critical intermediate computations (such as softmax normalization) are performed in FP32 precision. Furthermore, we find that FP8 \textsc{FlashAttention-3}—when using block quantization and incoherent execution—achieves 2.6$\times$ higher accuracy than standard attention leveraging per-tensor quantization, especially in situations featuring outlier values.

We have released \textsc{FlashAttention-3} under a permissive license\footnote{\textsc{FlashAttention-3} is available at \texttt{https://github.com/Dao-AILab/flash-attention}}, and aim to incorporate it into both PyTorch and Hugging Face libraries to maximize accessibility for researchers and developers.

\section{Background: Multi-Head Attention and GPU Characteristics}
\label{sec:background}

\subsection{Multi-Head Attention}
\label{subsec:multi_head_attn}

Let $\mathbf{Q}, \mathbf{K}, \mathbf{V} \in \mathbb{R}^{N \times d}$ be the query, key and value input sequences associated to a single head, where $N$ is the sequence length and $d$ is the head dimension. Then the attention output $\mathbf{O}$ is computed as:

\begin{equation*}
  \mathbf{S} = \alpha \mathbf{Q} \mathbf{K}^\top \in \mathbb{R}^{N \times N}, \quad \mathbf{P} = \mathrm{softmax}(\mathbf{S}) \in \mathbb{R}^{N \times N}, \quad \mathbf{O} = \mathbf{P}\mathbf{V} \in \mathbb{R}^{N \times d},
\end{equation*}

In this formulation, $\mathrm{softmax}$ operates along the rows, with the scaling parameter commonly chosen as $\alpha = 1/\sqrt{d}$. To address potential numerical issues caused by large exponentials, $\mathrm{rowmax}(\mathbf{S})$ is subtracted from $\mathbf{S}$ prior to applying the softmax. Within multi-head attention (MHA), individual heads utilize distinct query, key, and value projection matrices, and these calculations are conducted in parallel over various heads and across batches, generating the complete output tensor.

Now let $\phi$ be a scalar loss function and let $\mathbf{d}(-) = \partial \phi / \partial (-)$ be notation for the gradient. Given the output gradient $\mathbf{dO} \in \mathbb{R}^{N \times d}$, we compute $\mathbf{dQ}$, $\mathbf{dK}$, and $\mathbf{dV}$ according to the chain rule as follows:

\begin{align*}
  \mathbf{dV} &= \mathbf{P}^\top \mathbf{dO} \in \mathbb{R}^{N \times d} \\
  \mathbf{dP} &= \mathbf{dO} \mathbf{V}^\top \in \mathbb{R}^{N \times N} \\
  \mathbf{dS} &= \mathrm{dsoftmax} (\mathbf{dP}) \in \mathbb{R}^{N \times N} \\
  \mathbf{dQ} &= \alpha \mathbf{dS} \mathbf{K} \in \mathbb{R}^{N \times d} \\
  \mathbf{dK} &= \alpha \mathbf{dS}^\top \mathbf{Q} \in \mathbb{R}^{N \times d},
\end{align*}

In this context, we observe that $\mathbf{d}s = (\mathrm{diag}(p) - p p^\top)\mathbf{d}p$, where $p$ is given by $p = \mathrm{softmax}(s)$, viewing $p$ as a function of the input vector $s$. We denote by $\mathrm{dsoftmax}(\mathbf{dP})$ the application of this relationship across each row individually. This operation, like others in the backward propagation of MHA, naturally lends itself to parallel execution across multiple heads and batches.

\subsection{GPU hardware characteristics and execution model}
\label{subsec:hardware}

We discuss the key features of the GPU execution model pertinent to \textsc{FlashAttention-3}, concentrating specifically on the NVIDIA Hopper architecture to provide a concrete example of this framework.

\paragraph{Memory hierarchy:} On GPUs, memory is structured as a series of data domains, where smaller capacities are compensated by higher bandwidths, as outlined in \cref{tab:gpu-hierarchy}\footnote{\citet{luo2024benchmarking} details a shared memory bandwidth of 128 bytes per clock cycle for each SM; with 132 SMs operating at a boost clock of 1830 MHz, this cumulative bandwidth is obtained via multiplication.}. The global memory (GMEM), referred to as HBM, consists of off-chip DRAM and is universally accessible by all streaming multiprocessors (SMs). Information stored in GMEM is automatically cached within an on-chip L2 cache. Furthermore, each SM is equipped with its own compact, on-chip shared memory (SMEM)—an explicitly managed and extensively banked cache under programmer control. The memory hierarchy culminates in the register files, found inside each SM.

\paragraph{Thread hierarchy:} Within the GPU programming paradigm, execution units are logically arranged into several hierarchical layers. At the most granular level are threads, which aggregate into warps of 32 threads each. Next, four consecutive warps form a warpgroups. These are further grouped into threadblocks—also referred to as cooperative thread arrays (CTAs). In architectures such as Hopper, multiple threadblocks can be organized into threadblock clusters. At the top level, collections of threadblocks constitute grids.

There is a strong correlation between these two hierarchical structures. All threads belonging to a single CTA are assigned simultaneously to the same SM, while CTAs situated in a given cluster are jointly dispatched to the same GPC. The SMEM can be accessed by any thread inside the CTA, while RMEM, composed of up to 256 registers, is exclusively accessible to each individual thread.

\paragraph{Asynchrony and warp-specialization:}

GPUs function as throughput-oriented processors, leveraging concurrent and asynchronous operations to obscure delays associated with memory access and instruction execution. To support asynchronous memory transfers between GMEM and SMEM, the Hopper architecture introduces the Tensor Memory Accelerator (TMA), a specialized hardware mechanism \cite[\S7.29]{cuda}. Additionally, diverging from earlier designs like Ampere, Hopper's Tensor Core—accessible through the warpgroup-wide WGMMA instruction \cite[\S9.7.14]{ptx}—operates asynchronously and is capable of fetching its operands directly from shared memory.

Enabling hardware-level asynchrony facilitates the design of warp-specialized kernels, wherein the warps within a CTA are assigned exclusively to producer or consumer functions, performing either data transfer or computation, but not both. This separation generally enhances the compiler’s capability to construct more efficient instruction schedules \citep{warp-specialization-2011}. Moreover, Hopper introduces dynamic register reallocation between warpgroups using the \verb|setmaxnreg| instruction \cite[\S9.7.17.1]{ptx}. This mechanism allows warps engaged in MMA operations to access a greater portion of RMEM, in contrast to those that only perform TMA operations, which require the involvement of just a single thread.

\paragraph{Low-precision number formats:}
\label{sec:low-precision-gpu}
Recent advancements in GPU design include dedicated hardware modules to enhance low-precision computational tasks. As an illustration, the WGMMA instruction is capable of leveraging the FP8 Tensor Cores found in Hopper GPUs, achieving up to double the throughput per SM relative to FP16 or BF16.

Nevertheless, to properly utilize FP8 WGMMA, one must be aware of the specific operand layout requirements. Considering a GEMM operation involving $A \times B^{\top}$, where $A$ is of size $M\times K$ and $B$ is $N\times K$, we describe the $A$ or $B$ operand as \emph{mn-major} if its memory layout is contiguous with respect to the $M$ or $N$ (outer) dimension, whereas an operand is referred to as \emph{k-major} if the data is contiguous along the inner $K$ dimension. FP16 WGMMA accommodates both mn-major and k-major formats for input operands residing in SMEM. In contrast, FP8 WGMMA restricts supported operand layouts to the k-major form only. This restriction becomes particularly relevant in use cases such as attention mechanisms, where back-to-back GEMMs may be fused within a single kernel. Here, the incompatibility between the FP32 accumulator and the required FP8 operand layouts presents a challenge, complicating the invocation of sequential dependent FP8 WGMMAs.

When applied to attention mechanisms, these constraints on layout necessitate particular adjustments in the structure of an FP8 algorithm; a detailed discussion of these changes is presented in \cref{sec:algofp8}.

\subsection{Standard Attention and Flash Attention} As outlined by \citet{dao2022flashattention}, we define \textbf{standard attention} as the approach where attention computation on a GPU generates and stores the intermediate matrices $\mathbf{S}$ and $\mathbf{P}$ in HBM. The primary innovation of \textsc{FlashAttention} is the adoption of a local softmax reduction, which eliminates costly intermediate memory accesses by combining the attention process into one kernel execution. This local softmax operation occurs in lines \ref{code-ws:softmax_start}-\ref{code-ws:softmax_end} of the consumer mainloop within \cref{alg:flash3_wgmma_ws_only}, as well as through the normalization of $\mathbf{O}$’s sub-blocks. A concise proof that this method accurately computes $\mathbf{O}$ is provided in \cite[\S 2.3.1]{dao2023flashattention2}.

\section{FlashAttention-3: Algorithm}
\label{sec:algo}

This section provides an overview of the \textsc{FlashAttention-3} algorithm. Our exposition concentrates on the forward pass, while details regarding the backward pass can be found in~\cref{sec:algo_ws_bwd}. We begin by outlining the incorporation of warp-specialization and the use of a circular SMEM buffer within the foundational \textsc{FlashAttention-2} framework. Next, we discuss leveraging WGMMA’s asynchronous capabilities to establish a 2-stage pipeline that overlaps the GEMM and softmax computations. Lastly, we detail the required changes for handling FP8, addressing both layout compatibility and enhanced accuracy through block quantization and managing incoherent computation.

\subsection{Producer-Consumer asynchrony through warp-specialization and pingpong scheduling}
\label{sec:algo_ws}

\paragraph{Warp-specialization}
Similar to \textsc{FlashAttention-2}, the forward computation in \textsc{FlashAttention-3} exhibits a high degree of parallelism across the batch dimension, the number of attention heads, and the length of the query sequence. For this reason, we can focus our analysis at the CTA (cooperative thread array) level, where the algorithm processes a single tile $\mathbf{Q}_i$ from the query matrix to generate the matching output tile $\mathbf{O}_i$. For clarity, we initially present the warp-specialization approach utilizing a circular SMEM buffer, but \emph{without} incorporating the overlap between GEMM and softmax operations. Let $d$ denote the dimension of each head, $N$ the length of the sequence, and let us choose a block size $B_r$ for the queries so that $\mathbf{Q}$ is partitioned into $T_r = \lceil \frac{N}{B_r} \rceil$ blocks, specifically $\mathbf{Q}_1, .., \mathbf{Q}_{T_r}$.

\begin{algorithm}[H]
    \caption{\small\label{alg:flash3_wgmma_ws_only}\textsc{FlashAttention-3} forward pass \textbf{without} intra-consumer overlapping -- CTA view}
    \begin{algorithmic}[1]
\REQUIRE Matrices $\mathbf{Q}_i \in \mathbb{R}^{B_r \times d}$ as well as $\mathbf{K}, \mathbf{V} \in \mathbb{R}^{N \times d}$ located in HBM, key block size parameter $B_c$, and $T_c = \lceil \frac{N}{B_c} \rceil$.
\STATE Set up the pipeline construct to coordinate barrier synchronization, utilizing a circular SMEM buffer spanning $s$ stages.
\IF {running in producer warpgroup}
\STATE Free a specified amount of registers in advance.
\STATE Execute the load of $\mathbf{Q}_i$ from HBM to shared memory.
\STATE Upon the successful load, notify the consumer process that $\mathbf{Q}_i$ is now available.
\FOR{$0 \le j < T_c$}
    \STATE Hold until the $(j\,\%\,s)$th buffer stage is acknowledged as consumed.
    \STATE Initiate transfers of $\mathbf{K}_j, \mathbf{V}_j$ from HBM to shared memory targeting the $(j\,\%\,s)$th buffer stage.
    \STATE Once finished, signal consumers that the loading of $\mathbf{K}_j, \mathbf{V}_j$ has been completed.
\ENDFOR
\ELSE \STATE Allocate a specific number of registers anew, determined by how many consumer warps are utilized.
\STATE In on-chip memory, initialize $\mathbf{O}_i = (0) \in \mathbb{R}^{B_r \times d}$, $\ell_i = (0) \in \mathbb{R}^{B_r}$, and $m_i = (-\infty) \in \mathbb{R}^{B_r}$.
\STATE Wait for $\mathbf{Q}_i$ data to be present in shared memory.
\FOR{$0 \le j < T_c$}
\STATE Wait until $\mathbf{K}_j$ becomes accessible in shared memory.
\STATE Perform computation $\mathbf{S}_i^{(j)} = \mathbf{Q}_i \mathbf{K}_j^T$ (via SS-GEMM), then commit and stall until complete. \label{alg:ws_only_gemm1}
\STATE Save $m_i^{\mathrm{old}} = m_i$, then update $m_i = \mathrm{max}(m_i^{\mathrm{old}}, \mathrm{rowmax}(\mathbf{S}_i^{(j)}))$. \label{code-ws:softmax_start}
\STATE Calculate $\widetilde{\mathbf{P}}_i^{(j)} = \mathrm{exp}(\mathbf{S}_i^{(j)} - m_i)$ and update $\ell_i = \mathrm{exp}(m_i^{\mathrm{old}} - m_i) \ell_i + \mathrm{rowsum}(\widetilde{\mathbf{P}}_i^{(j)})$. \label{code-ws:softmax_end}
\STATE Await loading of $\mathbf{V}_j$ into shared memory.
\STATE Carry out $\mathbf{O}_i = \mathrm{diag}(\mathrm{exp}(m_i^{\mathrm{old}} - m_i))^{-1} \mathbf{O}_i + \widetilde{\mathbf{P}}_i^{(j)} \mathbf{V}_j$ (using RS-GEMM). Commit and synchronize. \label{alg:ws_only_gemm2}
\STATE Mark the $(j\,\%\,s)$th buffer stage as available for the producer process.
\ENDFOR
\STATE Finalize output via $\mathbf{O}_i = \mathrm{diag}(\ell_i)^{-1} \mathbf{O}_i$ and $L_i = m_i + \log(\ell_i)$.
\STATE Store $\mathbf{O}_i$ and $L_i$ to HBM, representing the $i$th partition of $\mathbf{O}$ and $L$.
\ENDIF
\end{algorithmic}
\end{algorithm}

In our realization of \cref{alg:flash3_wgmma_ws_only} on Hopper, we employ \verb|setmaxnreg| to manage (de)allocations, utilize TMA to load $\mathbf{Q}_i$ and $\{ \mathbf{K}_j, \mathbf{V}_j \}_{0 \leq j < T_c}$, and rely on WGMMA for executing the GEMMs within the main consumer loop; the choice between the SS or RS prefix distinguishes whether the primary operand resides in shared memory or the register file. When analyzing the execution order in \cref{alg:flash3_wgmma_ws_only}, it is important to recognize that TMA load operations function asynchronously and, therefore, do not block the issuance of subsequent loads. Additionally, during the first $s$ iterations of the producer’s mainloop, the buffer population process proceeds without inserting any waits.

\paragraph{Pingpong scheduling} Leveraging the asynchronous execution capabilities of WGMMA and TMA, coupled with warp specialization, enables the interleaving of softmax computations in one warpgroup with GEMM operations in another. This scheduling strategy is particularly appealing given that, on current hardware accelerators, non-matmul instructions tend to have substantially lower throughput than their matmul counterparts. For instance, on the H100 SXM5 GPU, FP16 matrix multiplication achieves 989 TFLOPS, whereas specialized functions like the exponential—crucial for softmax—achieve only 3.9 TFLOPS\footnote{According to the CUDA programming guide, each streaming multiprocessor (SM) can execute 16 special function operations per clock cycle; thus, the total throughput is estimated by multiplying 16 with 132 SMs and a 1830 MHz clock, resulting in 3.9 TFLOPS for special functions.}. In a typical FP16 attention forward pass with a head dimension of 128, the number of matmul FLOPs is 512 times that of the exponential operations, yet the latter are delivered at a throughput 256 times lower. As a result, computing exponentials can consume up to half the cycles required for matrix multiplication. This throughput disparity becomes even more pronounced in FP8 computations, where matmul performance doubles but the special function throughput remains fixed.

Given that the exponential function is executed by a distinct hardware component known as the multi-function unit, it is preferable to align the computation of exponentials concurrently with the matmul operations on the Tensor Cores. To orchestrate this overlap, we implement synchronization barriers (\texttt{bar.sync} instructions), which ensure that the GEMMs—specifically GEMM1 (\(\mathbf{P}\mathbf{V}\) from the current iteration) and GEMM0 (\(\mathbf{Q} \mathbf{K}^\top\) from the subsequent iteration)—in warpgroup 1 are dispatched ahead of those in warpgroup 2. This sequencing ensures that as warpgroup 2 undertakes its GEMM calculations, the softmax for warpgroup 1 can be processed in parallel. Afterwards, the work assignments switch: warpgroup 2 handles softmax while warpgroup 1 resumes GEMMs, resulting in a “pingpong” scheduling pattern, as depicted in~\cref{fig:pingpong_scheduling}. Although the actual runtime scheduling does not always conform precisely to the schematic representation, this technique generally yields noticeable performance improvements, such as increasing throughput from 570 TFLOPS to between 620 and 640 TFLOPS for FP16 forward passes with a head dimension of 128 and a sequence length of 8192.

\paragraph{Attention variants} In the case of multi-query attention~\citep{shazeer2019fast} and grouped query attention~\citep{ainslie2023gqa}, we adopt the methodology outlined in \textsc{FlashAttention-2}, modifying tensor indexing so as to prevent redundant copies of $\mathbf{K}$ and $\mathbf{V}$ within HBM.

\subsection{Intra-warpgroup overlapping GEMMs and softmax}

It is possible to achieve partial overlap between instructions related to the softmax operation and those involved in the GEMMs, even when restricted to a single warpgroup. One approach for facilitating this kind of overlap is outlined below.

In the attention algorithm, operations within the inner loop (main loop) have sequential dependencies that impede parallelization within a single iteration. For example, (local) softmax (lines \ref{code-ws:softmax_start} to \ref{code-ws:softmax_end}) relies on the output $\mathbf{S}_i^{(j)}$ of the first GEMM, while the second GEMM takes its result $\widetilde{\mathbf{P}}_i^{(j)}$ as an operand. Indeed, the wait statements in lines \ref{alg:ws_only_gemm1} and \ref{alg:ws_only_gemm2} of \cref{alg:flash3_wgmma_ws_only} serialize the execution of softmax and GEMMs. However, we can break these dependencies by pipelining across iterations through additional buffers in registers. Pursuing this idea, we propose the following two-stage\footnote{Note that the number of stages of the overlapping scheme is bounded by, but need not equal, the number $s$ of stages in the circular SMEM buffer.} GEMM-softmax pipelining algorithm:

\begin{algorithm}[H]
    \caption{\small\label{alg:flash3_wgmma}\textsc{FlashAttention-3} consumer warpgroup forward pass}
    \begin{algorithmic}[1]
      \REQUIRE  Matrices $\mathbf{Q}_i \in \mathbb{R}^{B_r \times d}$ and $\mathbf{K}, \mathbf{V} \in \mathbb{R}^{N \times d}$ located in HBM, along with key block size $B_c$ and $T_c = \lceil \frac{N}{B_c} \rceil$.
      \STATE Dynamically allocate a specified set of registers according to the number of consumer warps.
      \STATE On chip, set $\mathbf{O}_i = (0) \in \mathbb{R}^{B_r \times d}$, and assign $\ell_i = 0$, $m_i = -\infty \in \mathbb{R}^{B_r}$.
      \STATE Wait until both $\mathbf{Q}_i$ and the first key block $\mathbf{K}_0$ are present in shared memory.
      \STATE Calculate $\mathbf{S}_{\mathrm{cur}} = \mathbf{Q}_i \mathbf{K}_0^T$ utilizing WGMMA, then commit and synchronize.
      \STATE Free the stage 0 buffer slot dedicated to $\mathbf{K}$.
      \STATE Derive $m_{i}$, $\tilde{\mathbf{P}}_{\mathrm{cur}}$, and $\ell_{i}$ from $\mathbf{S}_{\mathrm{cur}}$, followed by rescaling $\mathbf{O}_i$ accordingly.
      \FOR{$1 \le j < T_c - 1$}
        \STATE Await transfer of key block $\mathbf{K}_j$ into shared memory.
        \label{code:mainloop_start}
        \STATE With WGMMA, compute $\mathbf{S}_{\mathrm{next}} = \mathbf{Q}_i \mathbf{K}_{j}^T$; commit without blocking. \label{code:first_wgmma}
        \STATE Await loading of value block $\mathbf{V}_{j-1}$ into shared memory.
        \STATE Update $\mathbf{O}_{i} = \mathbf{O}_{i} + \tilde{\mathbf{P}}_{\mathrm{cur}} \mathbf{V}_{j-1}$ by applying WGMMA; again, commit but continue execution. \label{code:second_wgmma}
        \STATE Wait for the result of the WGMMA product $\mathbf{Q}_i \mathbf{K}_{j}^T$.
        \STATE Update $m_{i}$, $\tilde{\mathbf{P}}_{\mathrm{next}}$, and $\ell_{i}$ via computations on $\mathbf{S}_{\mathrm{next}}$. \label{code:softmax}
        \STATE After the WGMMA $\tilde{\mathbf{P}}_{\mathrm{cur}} \mathbf{V}_{j-1}$ completes, apply the necessary scaling to $\mathbf{O}_i$.
        \STATE Release buffer stage $(j\,\%\,s)$ for $\mathbf{K}$, and $(j-1\,\%\,s)$ for $\mathbf{V}$.
        \STATE Assign $\mathbf{S}_{\mathrm{cur}}$ the content of $\mathbf{S}_{\mathrm{next}}$.
        \label{code:mainloop_end}
      \ENDFOR \STATE Await the arrival of value block $\mathbf{V}_{T_c - 1}$ in shared memory.
      \STATE Complete the computation with
      $\mathbf{O}_{i} = \mathbf{O}_{i} + \tilde{\mathbf{P}}_{\mathrm{last}} \mathbf{V}_{T_c - 1}$ using WGMMA, ensuring to commit and synchronize at this point.

\STATE Epilogue: Adjust $\mathbf{O}_{i}$ according to $m_{i}$. Derive $L_{i}$ utilizing both $m_{i}$ and $\ell_{i}$.
Store $\mathbf{O}_{i}$ and $L_{i}$ in HBM as the $i$-th segments of $\mathbf{O}$ and $L$, respectively.
\end{algorithmic}
\end{algorithm}

\cref{alg:flash3_wgmma} serves to replace the consumer pathway in \cref{alg:flash3_wgmma_ws_only}, thereby forming the full \textsc{FlashAttention-3} procedure for FP16 arithmetic. In general, WGMMA is used here to denote asynchronous GEMM. During the core loop (from lines \ref{code:mainloop_start} to \ref{code:mainloop_end}), the execution of the second WGMMA in step $j$ (line \ref{code:second_wgmma}) is concurrently performed with the softmax computation of the subsequent step, $j+1$ (line \ref{code:softmax}).

Although the pipelined scheme depicted earlier provides possible performance improvements in theory, several real-world limitations must be acknowledged:
\paragraph{Compiler reordering} While the pseudocode assumes a straightforward sequence of operations, in practice, the compiler (NVCC) may alter instruction order to achieve better optimizations. Such reordering may compromise the meticulously arranged alternation between WGMMA and non-WGMMA instructions, leading to suboptimal behavior or reduction of expected speedups. Examination of the SASS reveals, however, that the compiler typically produces the desired instruction interleaving (Section~\ref{sec:2-stage-sass}).
\paragraph{Register pressure} Achieving high performance depends on limiting register spilling, but the 2-stage pipeline design introduces increased register usage to hold intermediate computations and stage-specific context. In particular, an additional $\mathbf{S}_{\mathrm{next}}$ buffer needs to be allocated in registers, resulting in an overhead of $B_r \times B_c \times \text{sizeof}(\text{float})$ per threadblock. This elevation in register requirements can be at odds with other optimizations, such as enlarging block sizes, which likewise consume substantial register resources. Therefore, decisions regarding these trade-offs should be informed by profiling data.
\paragraph{3-stage pipelining} Beyond the 2-stage method above, we also consider a 3-stage pipeline where the second WGMMA operation and the softmax calculation are further interleaved. While this configuration could, in principle, enhance utilization of the Tensor Cores, it also intensifies register demands, making it more challenging to select appropriate tile sizes relative to pipeline depth. The specifics of the 3-stage implementation, along with performance analysis, appear in ~\cref{sec:3-stage}.

\subsection{Low-precision with FP8}
\label{sec:algofp8}

\textbf{Efficiency: layout transformations.} Executing the forward computation of \textsc{FlashAttention-3} using FP8 precision introduces unique challenges related to layout compatibility that do not arise when working with FP16.

To begin, it is important to observe that the input tensors $\mathbf{Q}$, $\mathbf{K}$, and $\mathbf{V}$ are conventionally arranged so that the head dimension is contiguous. However, for the second GEMM operation using FP8 WGMMA, which enforces the k-major constraint, we require $\mathbf{V}$—specifically, its tiles as loaded into SMEM—to be contiguous with respect to the sequence length dimension. Because the TMA load operation is incapable of altering which dimension is contiguous, there are two alternatives: (1) apply a transpose to $\mathbf{V}$ in GMEM during a pre-processing stage, or (2) perform an in-kernel transpose of tiles of $\mathbf{V}$ once they are in SMEM. In order to achieve (1), this can be approached by (1a) integrating the transpose with the epilogue of an earlier computation (e.g., rotary embedding), or (1b) invoking a dedicated pre-processing transpose kernel\footnote{A highly tuned transpose kernel can attain performance close to the device’s memory bandwidth~\citep{colfax_cutlass_transpose_2024}.}, which switches the strides of the sequence length and head axes. Nevertheless, option (1a) presents challenges when aiming for compatibility with standard libraries, while option (1b) leads to excessive overhead, particularly when memory bandwidth is the primary constraint, as in inference scenarios.

For FP8 \textsc{FlashAttention-3}, we select approach (2). The in-kernel transpose leverages the LDSM (\verb|ldmatrix|) and STSM (\verb|stmatrix|) instructions, enabling a warp of threads to collectively transfer data between SMEM and RMEM in units of 128 bytes.\footnote{According to the PTX documentation, LDSM/STSM instructions operate on $8 \times 8$ matrices with 16-bit values \cite[\S 9.7.13.4.15-16]{ptx}. To adapt these instructions for FP8, we pack pairs of 8-bit elements, permitting their use in this reduced precision scenario. However, the transpose variants of LDSM/STSM do not separate the packed 8-bit values, which makes additional register manipulation necessary between the LDSM and STSM stages to properly realize a tile-wise transpose; further specifics are omitted.} Because LDSM and STSM are efficient with respect to register usage, they can be invoked directly from within the producer warpgroup, while also providing flexibility to transpose memory layouts during the data movement. Additionally, subsequent to the initial iteration, it becomes possible to schedule the transposition of the next $\mathbf{V}$ tile concurrently with the execution of the two WGMMAs that operate on the previous $\mathbf{V}$ and current $\mathbf{K}$ tiles.

Second, we note that, in contrast to FP16, the FP32 accumulator’s memory layout in FP8 WGMMA operations does not match that of operand A when A resides in registers. Illustrative examples of these layouts are shown in \cref{fig:rmem_accum} and \cref{fig:rmem_operand}, which indicate the ordering of register-held entries per thread. Employing byte permute instructions enables us to rearrange the accumulator output from the initial WGMMA into a structure compatible both with the next WGMMA and with the $\mathbf{V}$ tile layout generated by the in-kernel transpose. More precisely, following the configuration depicted in \cref{fig:rmem_accum}, we permute the registers in the sequence:
$$\{ \verb|d0 d1 d4 d5 d2 d3 d6 d7| \},$$
and extend this permutation across each successive 8-byte block. This transformation permutes the columns of the logical $\mathbf{P}$ tile (such that, for example, columns $0189$ are shifted to the initial columns). To ensure WGMMA produces the intended output tile, the in-kernel transpose can, in turn, emit a corresponding row permutation for the $\mathbf{V}$ tile.\footnote{By leveraging this flexibility within the in-kernel transpose, there is no need to use shuffle instructions for redistributing register contents among threads, as was required in earlier approaches~\citep{colfax_fp8_flashattention_2024}.}

\textbf{Accuracy: block quantization and incoherent processing.} The FP8 (e4m3) numerical format allocates 3 bits for the mantissa and 4 bits to represent the exponent, resulting in greater quantization error in comparison to FP16 or BF16 formats. Furthermore, contemporary large-scale models tend to contain outlier values~\citep{dettmers2208llm, sun2024massive}, which exhibit magnitudes significantly exceeding those of most elements, posing additional challenges for quantization. Standard practice involves the application of per-tensor scaling~\citep{micikevicius2022fp8}, whereby each tensor is associated with a single scaling factor (for example, separate scalars for $\mathbf{Q}$, $\mathbf{K}$, and $\mathbf{V}$). 

To address and mitigate the elevated numerical errors observed in FP8-based attention computations, we introduce two strategies:

\begin{enumerate}

\item \textbf{Block quantization}: In this approach, a single scalar is retained per block, such that for each of $\mathbf{Q}$, $\mathbf{K}$, and $\mathbf{V}$, the tensor is partitioned into blocks of either $B_r \times d$ or $B_c \times d$, with quantization applied independently to each block. This quantization step can be integrated with operations that directly precede the attention mechanism (such as rotary embedding) without any added overhead, as these operations are typically constrained by memory bandwidth. Since the \textsc{FlashAttention-3} algorithm natively processes data in block format, each block of $\mathbf{S}$ can be appropriately scaled to accommodate the effects of block quantization with no additional computational cost.
\item \textbf{Incoherent processing}: To mitigate the influence of outliers, we pre-multiply both $\mathbf{Q}$ and $\mathbf{K}$ with a randomly chosen orthogonal matrix $\mathbf{M}$ prior to applying FP8 quantization. Because $\mathbf{M}$ is orthogonal, it holds that $\mathbf{M} \mathbf{M}^\top = I$, implying $(\mathbf{Q} \mathbf{M}) (\mathbf{K} \mathbf{M})^\top = \mathbf{Q} \mathbf{K}^\top$, and thereby ensuring the attention computation remains unchanged by this transformation. The effect of $\mathbf{M}$ is to distribute outliers across multiple entries, as each element in $\mathbf{Q} \mathbf{M}$ or $\mathbf{K} \mathbf{M}$ becomes a random linear combination of the respective entries in $\mathbf{Q}$ or $\mathbf{K}$, which decreases quantization errors. Following the methodology in \citet{chee2024quip} and~\citet{tseng2024quip}, $\mathbf{M}$ is constructed as a product of random diagonal matrices with entries $\pm 1$ and a Hadamard matrix. This formulation allows for multiplication to be performed in $O(d \log d)$ time complexity, rather than $O(d^2)$, and can similarly be integrated with the rotary embedding step without incurring extra computation.
\end{enumerate}

We provide empirical evidence in \cref{sec:numerical_error} showing these methods decrease numerical error by up to 2.6$\times$.

\section{Empirical Validation}
\label{sec:exp}

Our implementation of \textsc{FlashAttention-3} builds upon fundamental operations from CUTLASS~\citep{Thakkar_CUTLASS_2023}, leveraging the WGMMA and TMA abstractions to assess both its performance and correctness.

\begin{itemize}

\item \textbf{Attention benchmarking.} We evaluate the performance of \textsc{FlashAttention-3} in terms of runtime over varying sequence lengths, directly comparing it to the standard PyTorch implementation, \textsc{FlashAttention-2}, the Triton-based \textsc{FlashAttention-2} (leveraging H100-specific hardware features), as well as a vendor-optimized \textsc{FlashAttention-2} from cuDNN for H100 GPUs. Our analysis demonstrates that \textsc{FlashAttention-3} achieves speedups of up to 2.0$\times$ relative to \textsc{FlashAttention-2}, and up to 1.5$\times$ compared to the Triton-based version. Furthermore, \textsc{FlashAttention-3} attains a throughput of as much as 740 TFLOPs/s, corresponding to 75\% of the H100 GPU’s theoretical peak TFLOPs/s.
\item \textbf{Ablation analysis.} Through controlled experiments, we show that the enhancements stemming from warp-specialization and GEMM-softmax pipelining are key factors driving the improved efficiency of \textsc{FlashAttention-3}.
\item \textbf{FP8 attention accuracy.} We provide evidence that the combined use of block quantization and incoherent computation substantially diminishes the numerical error in FP8 computations for \textsc{FlashAttention-3} by a factor of 2.6$\times$.
\end{itemize}

\subsection{Benchmarking Attention}
\label{subsec:benchmark_attn}

We evaluate the performance of various attention algorithms by recording their runtimes on an H100 80GB SXM5 GPU across several scenarios, including configurations with or without causal masking and head dimensions set to 64 or 128, all using FP16 inputs. The outcomes are presented in~\cref{fig:benchmark_attn_fwd} and~\cref{fig:benchmark_attn_bwd}. Our benchmarks reveal that \textsc{FlashAttention-3} achieves approximately 1.5-2.0$\times$ greater speed in the forward pass and 1.5-1.75$\times$ in the backward pass relative to \textsc{FlashAttention-2}. When compared with conventional attention routines, \textsc{FlashAttention-3} demonstrates a remarkable speedup of up to 3-16$\times$. For input sequences of moderate to large size (exceeding 1k tokens), \textsc{FlashAttention-3} is able to outperform even specialized proprietary libraries, such as the closed-source cuDNN library optimized for H100 hardware.

\paragraph{Benchmark settings:} We vary the sequence length as 512, 1k, ..., 16k, and set batch size so that the total number of tokens is 16k. We set the hidden dimension to 2048, and head dimension to be either 64, 128, or 256 (i.e., 32 heads, 16 heads, or 8 heads). To calculate the FLOPs of the forward pass, we use:

\begin{equation*}
  4 \cdot \text{seqlen}^2 \cdot \text{head dimension} \cdot \text{number of heads}.
\end{equation*}

Applying causal masking, the value is halved because roughly 50\% of the entries are actually computed. For the backward pass, we estimate FLOPs by multiplying the forward pass FLOPs by 2.5, reflecting the presence of 2 matrix multiplications during the forward computation and 5 during the backward computation as a result of recomputation.

Additionally, we evaluate the forward pass runtime for FP8 using comparable conditions. The performance outcomes for a head dimension of 256 are shown in ~\cref{fig:benchmark_attn_fp8}, while comprehensive results are presented in \cref{sec:benchmark_attn_fp8_full}.

\subsection{Ablation Study: 2-Stage Pipelining Experiments}
\label{sec:2-stage-pipelining-experiments}

We perform an ablation study on the two-stage WGMMA-softmax pipeline and on warp-specialization within the context of non-causal FP16 \textsc{FlashAttention-3}, maintaining fixed parameter settings $\{\text{batch}, \text{seqlen}, \text{nheads}, \text{hdim}\} = \{ 4, 8448, 16, 128\}$. As demonstrated in~\cref{table:ablation_pipelining}, our algorithmic enhancements—namely, introducing asynchrony through warp-specialization and enabling overlap between the GEMM and softmax computations—produce a marked performance improvement, increasing throughput from 570 to 661 TFLOPs.

\subsection{Numerical Error Validation}
\label{sec:numerical_error}

Given the recent attention on the numerical inaccuracies of \textsc{FlashAttention}~\citep{golden2024flash}, we conduct a comparison between \textsc{FlashAttention-2}, \textsc{FlashAttention-3}, and a typical baseline attention implementation, using an FP64 reference as a standard. In order to model the presence of outlier features and activations commonly observed in LLMs~\citep{dettmers2208llm, sun2024massive}, we initialize the elements of $\mathbf{Q}, \mathbf{K}, \mathbf{V}$ based on the following probability distribution:

\begin{equation*}
  \mathcal{N}(0, 1) + \mathcal{N}(0, 100) \cdot \mathrm{Bernoulli}(0.001).
\end{equation*}

In this setup, individual entries follow a normal distribution with a mean of zero and a standard deviation of 1; however, for 0.1\% of these entries, we introduce an additional independent normal variable with a standard deviation of 10. Subsequently, we compute the root mean squared error (RMSE) as displayed in~\cref{table:numerical_error}. When using FP16, both \textsc{FlashAttention-2} and \textsc{FlashAttention-3} yield an RMSE that is 1.7$\times$ lower than that of the conventional implementation, as they retain intermediate results (softmax) in FP32. For FP8, the baseline attention approach employs per-tensor scaling, with matmul accumulation performed in FP32 and intermediate softmax outputs maintained in FP16. Due to the incorporation of block quantization and non-coherent computation, \textsc{FlashAttention-3} operating in FP8 achieves 2.6$\times$ greater accuracy relative to this baseline.

\section{Dicussion, Limitations, Conclusion}
\label{sec:discussion}

Through the development of \textsc{FlashAttention-3}, we illustrate how leveraging novel programming methods alongside advanced hardware capabilities, including asynchronous execution and low-precision computation, can substantially enhance both the efficiency and fidelity of attention mechanisms. Our approach delivers a 1.5-2.0$\times$ speedup relative to \textsc{FlashAttention-2}, and achieves a 2.6$\times$ reduction in FP8 numerical errors compared to conventional per-tensor quantization approaches. Nevertheless, there are areas that warrant further investigation: improving optimization for LLM inference, incorporating persistent kernel strategies into the FP8 kernel,\footnote{In our experimental setup, FP16 \textsc{FlashAttention-3} benefits from a persistent kernel and a load balancing approach, whereas the FP8 version lacks these enhancements. This partially accounts for the inferior performance of FP8 \textsc{FlashAttention-3} at lower sequence lengths and under causal masking conditions when contrasted with FP8 cuDNN kernels.} and rigorously analyzing the implications of low-precision attention during large-scale training. While our results primarily target Hopper GPUs, we anticipate that the methods introduced will extend to various hardware accelerators. Ultimately, we aspire for this more rapid and accurate attention primitive to foster advancements and innovation in long-context tasks.

\end{document}